\section{Related Work}

\paragraph{Scaling-law forms and reliability.}
Empirical power laws have been reported across language, vision, and generative modeling \citep{hestness2017deep,rosenfeld2020constructive,kaplan2020scaling,henighan2020scaling,hoffmann2022empirical}. Smoothly broken laws model observed regime changes \citep{caballero2022broken}, while downstream studies show that simple trends are not universal \citep{ruan2024observational,lourie2025unreliable}. We do not propose another point-fit form; we characterize the future values that cannot be ruled out.

\paragraph{Spectral theory.}
Power-law eigenspectra and source conditions yield parameter--data--compute laws in kernel, random-feature, and linear models \citep{lin2024scaling,bordelon2024dynamical,lin2025improved}. Our population truncation model isolates the spectral tail exactly, permitting nonasymptotic testing constructions and a finite-sample OLS bridge.

\paragraph{Uncertainty and identifiability.}
Learning-curve extrapolation supports early stopping and hyperparameter optimization \citep{domhan2015speeding,klein2017fast}. Bayesian and prior-data fitted methods produce posterior predictions conditional on synthetic-function priors \citep{adriaensen2023efficient,lee2025bayesian}. Concurrent conformalized scaling laws calibrate residual-based prediction intervals \citep{romano2019conformalized,bukkapatnam2026conformalized}. \ScaleCert\ instead gives conditional frequentist coverage over an explicit spectral class and may return an unbounded interval when the pilot data do not identify the target. Kricheli \etal\ show local ill-conditioning when two-resource laws are fit along nearly collinear tokens-per-parameter rays \citep{kricheli2026coverage}. Our obstruction is one-dimensional, allows hidden future regimes, applies to arbitrary pilot placement inside a bounded interval, and is paired with target-specific partial identification. We use standard two-point testing arguments \citep{tsybakov2009introduction}.

\section{Problem Setup}

Let $M_1<\cdots<M_m\le M_{\max}$ be pilot model sizes and $M_\star>M_{\max}$ a target. We observe
\begin{equation}
  Y_i=\Risk(M_i)+\xi_i,\qquad \xi_i\sim\Normal(0,s_i^2),
  \label{eq:observations}
\end{equation}
independently. Standard errors may come from repeated training seeds or evaluation uncertainty. We distinguish the asymptotic exponent, the future risk $\Risk(M_\star)$, and the decision $\Risk(M_\star)\le\tau$; these need not have the same identifiability.

For $E_{\Risk}=\lim_{M\to\infty}\Risk(M)$, define the excess-risk exponent, when it exists, by
\begin{equation}
 a(\Risk):=-\lim_{M\to\infty}
 \frac{\log(\Risk(M)-E_{\Risk})}{\log M}.
 \label{eq:exponent-functional}
\end{equation}
For the one-power lower bound, write $t=\log(M/M_0)$. Pilots lie in $[-T,T]$ and the target is $t_\star=T+h$; $h/T$ is the extrapolation horizon relative to observed log-span.
